% Lumière du Soleil traversant l'atmosphère à midi et au coucher (énoncé du défi D19)
\begin{tikzpicture}[x=1cm,y=1cm]
\fill[otNavy!8] (-5.2,0) arc (180:0:5.2) -- (4.4,0) arc (0:180:4.4) -- cycle;
\fill[otGreen!30] (-4.4,0) arc (180:0:4.4) -- cycle;
\draw[otNavy!40,line width=0.5pt] (-5.2,0) arc (180:0:5.2);
\node[pcLegende,anchor=south,text=otNavy] at (0,5.25) {atmosphère (épaisseur très exagérée)};
\fill[otInk] (0,4.4) circle (0.07);
\node[pcLegende,anchor=north] at (0,4.32) {observateur};
\draw[orange,line width=1.2pt,-{Stealth[length=2.2mm]}] (0,6.0) -- (0,4.5) node[pcLegende,anchor=south west,pos=0.1,text=orange!80!black] {Soleil à midi};
\draw[orange,line width=1.2pt,-{Stealth[length=2.2mm]}] (-5.6,4.25) -- (-0.1,4.4);
\node[pcLegende,anchor=south,text=orange!80!black] at (-3.6,4.35) {Soleil au coucher : long trajet dans l'air};
\draw[otNavy,line width=0.6pt,-{Stealth[length=1.5mm]}] (-2.0,4.35) -- (-2.4,3.8);
\draw[otNavy,line width=0.6pt,-{Stealth[length=1.5mm]}] (-1.2,4.37) -- (-0.9,4.95);
\node[pcLegende,anchor=north,text=otNavy] at (-2.4,3.75) {bleu diffusé};
\end{tikzpicture}
\par\smallskip\legende{Les molécules de l'air diffusent la lumière dans toutes les directions, surtout les courtes longueurs d'onde.}
